\section{Discussion and Limitations}
\label{sec:discussion}

scVICAR uses intercellular structure to construct corrupted training examples.
The decoder recovers the original cell, and the encoder processes clean
expression profiles at inference. This design adds local information to scMAE
with a single graph operation during training. The convex-hull and displacement
bounds describe the resulting input change.

The full scVICAR pipeline performed well across the evaluated datasets. Among
16 primary configurations in the 15-dataset benchmark, scVICAR-T and scVICAR-F ranked
first and second on all four mean metrics. In the six frozen datasets,
scVICAR-T had the highest mean ARI, and scVICAR-F had the highest mean NMI and
macro-F1. The matched-backbone comparisons showed smaller
effects: $-0.002$ ARI for scVICAR-F and $+0.005$ for scVICAR-T relative to
NoMix. The external rankings describe complete pipelines, whereas the matched
effects estimate the incremental contribution of local mixing. That
contribution varied by dataset.

The ablations clarify the role of each component. RandomMix performed worse
than local mixing, indicating that neighbor selection matters. The gate-only
variant had the highest descriptive clean-graph ARI, whereas edge-only
weighting remained below NoMix. Cell-specific control of mixing strength was
the stronger adaptive component on the clean confirmatory graphs.
The full model also provides edge weights that can be examined under graph
corruption.

The stress experiment measured that response directly. Topology-informed
affinity separated retained same-class edges from injected cross-class edges
with mean AUROC 0.918. The cell-specific mixing coefficient correlated with
weighted neighborhood purity (mean Spearman 0.520). ARI degradation favored
scVICAR-T on two datasets and scVICAR-F on one at complete edge replacement.
Graph geometry and the encoder response therefore determine whether accurate
edge scoring improves the final clustering.

The downstream tasks showed a similarly specific pattern. Frozen low-label
probes improved by 0.13--0.34\% in relative macro-F1, with positive changes on
most datasets. Marker recovery and marker-overlap annotation varied more
widely, and scVICAR-T had negative mean changes on both marker endpoints.
Local mixing changed linear class separation more consistently than
cluster-level differential-expression programs. High baseline probe accuracy and
dataset-specific marker quality may also limit the observed gains.

Several design choices restrict the current method. The static PCA cosine KNN
graph can connect different cell types near transitions, batch boundaries, and
rare populations. The topology-informed affinity is an uncalibrated analytic
weight. Cosine similarity and cosine distance are algebraically redundant in
its current form. The sampler also uses the same affinity for sampling and
reweighting, which sharpens the effective kernel. Alternative estimators in the
ablation produced similar mean ARI at zero injected contamination.

The evaluation covers six frozen datasets, three model seeds, known-$K$
KMeans, and fixed-resolution Leiden. Other tissues, disease states, platforms,
and batch structures may produce different graph errors. The complete-label
sensitivity analysis suggests a possible benefit for rare or unresolved
populations and motivates a dedicated rare-cell evaluation. The low-label
probe measures within-dataset transductive annotation. Within this scope,
scVICAR provides competitive masked representations and explicit measurements
of how local graph quality affects neighbor mixing.
